\section{أنظمة التغذية الراجعة (\textenglish{Feedback Systems})}
\label{sec:feedback_systems}

السمة الجوهرية التي تُميز محرك السيرفو عن بقية المحركات الكهربائية التقليدية هي وجود حساس تغذية راجعة مُدمج، يُثبت عادةً في الجزء الخلفي من عمود الدوران (\textenglish{Motor Shaft}). تكمن وظيفة هذا الحساس في قراءة الموضع الزاوي، السرعة، واتجاه الدوران بشكل لحظي، ثم إرسال هذه البيانات إلى وحدة التحكم لإغلاق حلقة التحكم وتصحيح مسار الحركة.

تُقسم أنظمة التغذية الراجعة المستخدمة في الصناعة إلى عدة أنواع بناءً على تقنية القياس ومستوى الدقة المطلوب:

\subsection{المشفرات (\textenglish{Encoders})}
تُعد المشفرات (الضوئية أو المغناطيسية) الحساسات الأكثر شيوعاً في تطبيقات الأتمتة الدقيقة. وتعتمد على تحويل الحركة الميكانيكية إلى إشارات رقمية يفهمها المتحكم. تنقسم المشفرات إلى فئتين أساسيتين:

\begin{enumerate}
    \item \textbf{المشفرات التزايدية أو النسبية (\textenglish{Incremental Encoders}):}
    يقوم هذا النوع بتوليد سلسلة من النبضات الكهربائية (\textenglish{Pulses}) المتساوية مع كل زيادة في زاوية الدوران. يقوم المتحكم بحساب الموضع الفعلي عن طريق عَدّ هذه النبضات انطلاقاً من نقطة مرجعية. 
    \begin{itemize}
        \item \textbf{الميزة:} بساطة التصميم وانخفاض التكلفة.
        \item \textbf{العيب:} عند انقطاع التيار الكهربائي، يفقد النظام قراءة الموضع الحالي، مما يستدعي إجراء عملية "العودة إلى نقطة الصفر" (\textenglish{Homing Procedure}) عند إعادة التشغيل لتعريف الآلة بمكانها الصحيح.
    \end{itemize}
    
    \item \textbf{المشفرات المطلقة (\textenglish{Absolute Encoders}):}
    على عكس النوع التزايدي، يقرأ هذا المشفر الموقع باستخدام قرص مشفر يحمل نمطاً معقداً، بحيث يقابل كل موضع زاوي كوداً رقمياً فريداً (مثل \textenglish{Gray Code}).
    \begin{itemize}
        \item \textbf{الميزة الجوهرية:} يحتفظ بقيمة الموضع الفعلي والدقيق حتى في حال انقطاع الطاقة الكهربائية، مما يجعله الخيار القياسي والآمن في التطبيقات الحرجة كالأذرع الروبوتية، حيث قد تتسبب عملية الـ \textenglish{Homing} باصطدام الروبوت بالمعدات المحيطة.
    \end{itemize}
\end{enumerate}

\subsection{المحللات أو الريزولفرات (\textenglish{Resolvers})}
المحلل هو عبارة عن محول كهربائي دوار (\textenglish{Rotary Electrical Transformer}) يُستخدم لقياس زاوية الدوران. على عكس المشفرات التي تعطي إشارات رقمية، يقوم المحلل بتوليد إشارات تماثلية (\textenglish{Analog Signals}) متمثلة في موجات جيبية وجيب تمام (\textenglish{Sine and Cosine}) تتناسب سعتها مع زاوية العضو الدوار.

رغم أن المحللات قد توفر دقة أقل نسبياً من المشفرات الضوئية المتقدمة، إلا أنها تتفوق في جانب بالغ الأهمية: \textbf{المتانة الميكانيكية والاعتمادية}. لكونها تعتمد على ملفات نحاسية وقلب حديدي (بدون دوائر إلكترونية دقيقة أو أقراص زجاجية قابلة للكسر)، فإنها تتحمل ظروف العمل الصناعية القاسية جداً، مثل درجات الحرارة المرتفعة، والاهتزازات العنيفة، والأتربة، مما يجعلها مفضلة في الصناعات الثقيلة والتطبيقات العسكرية.